%=========================================================================
%  Préambule commun — Analyse des données, L3 Économie
%  Appelé par chaque séance via %=========================================================================
%  Préambule commun — Analyse des données, L3 Économie
%  Appelé par chaque séance via %=========================================================================
%  Préambule commun — Analyse des données, L3 Économie
%  Appelé par chaque séance via %=========================================================================
%  Préambule commun — Analyse des données, L3 Économie
%  Appelé par chaque séance via \input{preambule-cy}
%=========================================================================
\usetheme[progressbar=frametitle,numbering=fraction,block=fill]{metropolis}

\usepackage[T1]{fontenc}
\usepackage[utf8]{inputenc}
\usepackage[french]{babel}
%\usepackage{lmodern}  % absent de cet environnement de compilation ; a decommenter
\usepackage{booktabs}
\usepackage{listings}
\usepackage{xcolor}

%-------------------------------------------------------------------------
%  PALETTE — remplacer par les valeurs de votre charte CY habituelle
%-------------------------------------------------------------------------
\definecolor{CYprincipal}{HTML}{1B2A4A}   % couleur dominante (titres, barres)
\definecolor{CYaccent}{HTML}{C8102E}      % couleur d'accentuation
\definecolor{CYgris}{HTML}{55637C}        % texte secondaire
\definecolor{CYclair}{HTML}{EEF1F6}       % fonds de blocs

\setbeamercolor{palette primary}{bg=CYprincipal,fg=white}
\setbeamercolor{progress bar}{fg=CYaccent}
\setbeamercolor{frametitle}{bg=CYprincipal,fg=white}
\setbeamercolor{alerted text}{fg=CYaccent}
\setbeamercolor{block title}{bg=CYclair,fg=CYprincipal}
\setbeamercolor{block body}{bg=CYclair!60,fg=black}

%-------------------------------------------------------------------------
%  Mise en forme du code Python
%-------------------------------------------------------------------------
\lstdefinestyle{py}{
  language=Python,
  basicstyle=\ttfamily\small,
  keywordstyle=\color{CYaccent},
  stringstyle=\color{CYprincipal},
  commentstyle=\color{CYgris}\itshape,
  showstringspaces=false,
  breaklines=true,
  frame=none,
  literate={é}{{\'e}}1 {è}{{\`e}}1 {à}{{\`a}}1 {ç}{{\c c}}1 {ê}{{\^e}}1 {ô}{{\^o}}1 {û}{{\^u}}1 {î}{{\^i}}1
}
\lstset{style=py}

\author{Stefania Marcassa}
\institute{CY Cergy Paris Université --- L3 Économie}
\date{}

%=========================================================================
\usetheme[progressbar=frametitle,numbering=fraction,block=fill]{metropolis}

\usepackage[T1]{fontenc}
\usepackage[utf8]{inputenc}
\usepackage[french]{babel}
%\usepackage{lmodern}  % absent de cet environnement de compilation ; a decommenter
\usepackage{booktabs}
\usepackage{listings}
\usepackage{xcolor}

%-------------------------------------------------------------------------
%  PALETTE — remplacer par les valeurs de votre charte CY habituelle
%-------------------------------------------------------------------------
\definecolor{CYprincipal}{HTML}{1B2A4A}   % couleur dominante (titres, barres)
\definecolor{CYaccent}{HTML}{C8102E}      % couleur d'accentuation
\definecolor{CYgris}{HTML}{55637C}        % texte secondaire
\definecolor{CYclair}{HTML}{EEF1F6}       % fonds de blocs

\setbeamercolor{palette primary}{bg=CYprincipal,fg=white}
\setbeamercolor{progress bar}{fg=CYaccent}
\setbeamercolor{frametitle}{bg=CYprincipal,fg=white}
\setbeamercolor{alerted text}{fg=CYaccent}
\setbeamercolor{block title}{bg=CYclair,fg=CYprincipal}
\setbeamercolor{block body}{bg=CYclair!60,fg=black}

%-------------------------------------------------------------------------
%  Mise en forme du code Python
%-------------------------------------------------------------------------
\lstdefinestyle{py}{
  language=Python,
  basicstyle=\ttfamily\small,
  keywordstyle=\color{CYaccent},
  stringstyle=\color{CYprincipal},
  commentstyle=\color{CYgris}\itshape,
  showstringspaces=false,
  breaklines=true,
  frame=none,
  literate={é}{{\'e}}1 {è}{{\`e}}1 {à}{{\`a}}1 {ç}{{\c c}}1 {ê}{{\^e}}1 {ô}{{\^o}}1 {û}{{\^u}}1 {î}{{\^i}}1
}
\lstset{style=py}

\author{Stefania Marcassa}
\institute{CY Cergy Paris Université --- L3 Économie}
\date{}

%=========================================================================
\usetheme[progressbar=frametitle,numbering=fraction,block=fill]{metropolis}

\usepackage[T1]{fontenc}
\usepackage[utf8]{inputenc}
\usepackage[french]{babel}
%\usepackage{lmodern}  % absent de cet environnement de compilation ; a decommenter
\usepackage{booktabs}
\usepackage{listings}
\usepackage{xcolor}

%-------------------------------------------------------------------------
%  PALETTE — remplacer par les valeurs de votre charte CY habituelle
%-------------------------------------------------------------------------
\definecolor{CYprincipal}{HTML}{1B2A4A}   % couleur dominante (titres, barres)
\definecolor{CYaccent}{HTML}{C8102E}      % couleur d'accentuation
\definecolor{CYgris}{HTML}{55637C}        % texte secondaire
\definecolor{CYclair}{HTML}{EEF1F6}       % fonds de blocs

\setbeamercolor{palette primary}{bg=CYprincipal,fg=white}
\setbeamercolor{progress bar}{fg=CYaccent}
\setbeamercolor{frametitle}{bg=CYprincipal,fg=white}
\setbeamercolor{alerted text}{fg=CYaccent}
\setbeamercolor{block title}{bg=CYclair,fg=CYprincipal}
\setbeamercolor{block body}{bg=CYclair!60,fg=black}

%-------------------------------------------------------------------------
%  Mise en forme du code Python
%-------------------------------------------------------------------------
\lstdefinestyle{py}{
  language=Python,
  basicstyle=\ttfamily\small,
  keywordstyle=\color{CYaccent},
  stringstyle=\color{CYprincipal},
  commentstyle=\color{CYgris}\itshape,
  showstringspaces=false,
  breaklines=true,
  frame=none,
  literate={é}{{\'e}}1 {è}{{\`e}}1 {à}{{\`a}}1 {ç}{{\c c}}1 {ê}{{\^e}}1 {ô}{{\^o}}1 {û}{{\^u}}1 {î}{{\^i}}1
}
\lstset{style=py}

\author{Stefania Marcassa}
\institute{CY Cergy Paris Université --- L3 Économie}
\date{}

%=========================================================================
\usetheme[progressbar=frametitle,numbering=fraction,block=fill]{metropolis}

\usepackage[T1]{fontenc}
\usepackage[utf8]{inputenc}
\usepackage[french]{babel}
%\usepackage{lmodern}  % absent de cet environnement de compilation ; a decommenter
\usepackage{booktabs}
\usepackage{listings}
\usepackage{xcolor}

%-------------------------------------------------------------------------
%  PALETTE — remplacer par les valeurs de votre charte CY habituelle
%-------------------------------------------------------------------------
\definecolor{CYprincipal}{HTML}{1B2A4A}   % couleur dominante (titres, barres)
\definecolor{CYaccent}{HTML}{C8102E}      % couleur d'accentuation
\definecolor{CYgris}{HTML}{55637C}        % texte secondaire
\definecolor{CYclair}{HTML}{EEF1F6}       % fonds de blocs

\setbeamercolor{palette primary}{bg=CYprincipal,fg=white}
\setbeamercolor{progress bar}{fg=CYaccent}
\setbeamercolor{frametitle}{bg=CYprincipal,fg=white}
\setbeamercolor{alerted text}{fg=CYaccent}
\setbeamercolor{block title}{bg=CYclair,fg=CYprincipal}
\setbeamercolor{block body}{bg=CYclair!60,fg=black}

%-------------------------------------------------------------------------
%  Mise en forme du code Python
%-------------------------------------------------------------------------
\lstdefinestyle{py}{
  language=Python,
  basicstyle=\ttfamily\small,
  keywordstyle=\color{CYaccent},
  stringstyle=\color{CYprincipal},
  commentstyle=\color{CYgris}\itshape,
  showstringspaces=false,
  breaklines=true,
  frame=none,
  literate={é}{{\'e}}1 {è}{{\`e}}1 {à}{{\`a}}1 {ç}{{\c c}}1 {ê}{{\^e}}1 {ô}{{\^o}}1 {û}{{\^u}}1 {î}{{\^i}}1
}
\lstset{style=py}

\author{Stefania Marcassa}
\institute{CY Cergy Paris Université --- L3 Économie}
\date{}
